\newpage
\section{Lecture-12}
\subsection{Implicit Function Theorem}
Let's look $x^2+y^2=1$ closely. Is this a graph of a function? It's not.
No, it's not. It's a relation between $x$ and $y$. But we want to study function so far. 
At a neighborhood of $(x,y)$ except $(0,1)$ and $(0,-1)$, we can solve for $y$ as a function of $x$. 
$y=\sqrt{1-x^2}$ or $y=-\sqrt{1-x^2}$ are explicit functions that give the same relation between $x$ and $y$
as $x^2+y^2=1$. This is true locally near any point except $(0,1)$ and $(0,-1)$.

The implicit function theorem provides conditions under which a relationship of the from $F(x,y)=0$ can be rewritten as
a function $y=f(x)$ locally. 
\begin{example}
    Consider $y^5+y^3+y+x=0$. It is trivial to write $x=f(y)$, but can we write $y=f(x)$?
    But there is no explicit formula for $y$ in terms of $x$, because there is no formula for the roots of a general quintic polynomial.
    In general, there is no formula for finding roots of polynomials of degree 5 or higher (Abel, Galois). 
    Fix $x_0$, define $\varphi(y)=y^5+y^3+y+x_0$. Then,
\[
\varphi'(y)=5y^4+3y^2+1>0,\qquad\forall y\in\R.
\] 
Which implies $\varphi$ is strictly increasing. In Calculus-I, we know that any odd degree polynomial has at least one real root.
But since $\varphi$ is strictly increasing, it has exactly one real root. So we can write $y=f(x)$ locally near any point $(x_0,y_0)$, give me 
$x_0$, we can find a unique $y_0$ such that $y_0^5+y_0^3+y_0+x_0=0$. We found a function $y=f(x)$, but $f$ is implicit.       
\end{example} 
\begin{theorem}
Denote $\vec x=(x_1,\ldots,x_n)$. Let $F(\vec x,y)\in C^1$ in a neighborhood of $N_0(\vec x_0,y_0)$ such that
\begin{enumerate}
    \item $F(x_0,y_0)=0$,
    \item $\dfrac{\partial F}{\partial y}(x_0,y_0)\neq 0$.
\end{enumerate}
Then there exists a neighborhood of \((x_0,y_0)\) in which there is a unique (implicit) function
$y=f(\vec x)$ such that
\begin{enumerate}
    \item[(i)] \(y_0=f(\vec x_0)\),
    \item[(ii)] \(F\bigl(\vec x,f(\vec x)\bigr)=0\) for every \(x\) in the neighborhood,
    \item[(iii)] And the derivative of \(f\) is given by,
    \[
    \frac{\partial f}{\partial x_i}=-\frac{\dfrac{\partial F}{\partial x_i}}{\dfrac{\partial F}{\partial y}}
    \]
    in the neighborhood of $(\vec x_0,y_0)$.
\end{enumerate}
\end{theorem} 
\begin{example}
Given that $3x^2y - yz^2 - 4xz - 7 = 0$. We will show that near \((-1,1,2)\) we can write
$y = f(x,z)$, and we will find $\frac{\partial f}{\partial x}(-1,2)$.
Here,
\[
F(x,y,z) = 3x^2y - yz^2 - 4xz - 7 \in C^1
\]
And $F(-1,1,2) = 0$. Moreover,
\[
\frac{\partial F}{\partial y} = 3x^2 - z^2,
\qquad
\frac{\partial F}{\partial y}(-1,1,2) = -1 \ne 0.
\]
Thus, by the implicit function theorem, there exists a neighborhood of \((-1,1,2)\) in which there is a unique function \(y=f(x,z)\in C^1\) such
that \(f(-1,2)=1\) and \(F\bigl(x,f(x,z),z\bigr)=0\) for every \((x,z)\) in the neighborhood. Moreover,
\[
\frac{\partial f}{\partial x}(-1,2) = -\frac{\frac{\partial F}{\partial x}}{\frac{\partial F}{\partial y}}(-1,1,2) = -\frac{6xy - 4z}{3x^2 - z^2}(-1,1,2) = -\frac{6(-1)(1) - 4(2)}{3(-1)^2 - (2)^2} = -\frac{-6 - 8}{3 - 4} = -\frac{-14}{-1} = -14.
\]
In fact, 
\begin{align*}
    3x^2y-yz^2&= 4xz+7\\
    y(3x^2-z^2) &= 4xz+7\\
    y&=\frac{4xz+7}{3x^2-z^2}.
\end{align*}
In this example, we can explicitly solve for \(y\) in terms of \(x\) and \(z\), but the implicit function theorem guarantees the existence of such a function even when an explicit formula is not available. The key point is that the condition \(\frac{\partial F}{\partial y}(-1,1,2) \neq 0\) ensures that we can locally solve for \(y\) as a function of \(x\) and \(z\).
In this example, we can write $z=f(x,y)$ as well, but the theorem fails at \((-1,1,2)\) because
\[\frac{\partial F}{\partial z}(-1,1,2) = 0.\]
\end{example}
\begin{example}
\[
\text{ex.}\qquad F(x,y)=(x-y)^3
\]
\[
F(x,y)=0 \iff y=x
\]
There is an explicit function at any point.
However, at \((0,0)\),
\[
\frac{\partial F}{\partial x}(0,0)=3(x-y)^2\big|_{(0,0)}=0,
\]
\[
\frac{\partial F}{\partial y}(0,0)=-3(x-y)^2\big|_{(0,0)}=0.
\]
The theorem doesn't apply.
\end{example}
Let's prove the implicit function theorem for the case of two variables only.  
\begin{proof}
    $\frac{\partial F}{\partial y}(x_0,y_0)\neq 0$ implies w.l.o.g. $\frac{\partial F}{\partial y}(x_0,y_0)>0$.
    Since, $F\in C^1$, $\frac{\partial F}{\partial y}$ is continuous, so $\frac{\partial F}{\partial y}>0$ 
    at a neighborhood of $(x_0,y_0)$. What do we know about $F(x_0,y)$ function? Here we fixing $x_0$,
    and view $F$ as a function of $y$. We know that $\frac{\partial F}{\partial y}>0$ at a neighborhood of $(x_0,y_0)$, which implies
    $F(x_0,y)$ is a strictly increasing function. $\exists y_1$ sucht $F(x_0,y_1)>0$ and $\exists y_2$ such that $F(x_0,y_2)<0$.
    Since $F\in C^1$, $F$ is continuous. Which implies $\forall x$ near $x_0$ (releasing our first variable),
    \[ F(x,y_1)>0, \qquad F(x,y_2)<0\]
    For such an $x$ near $x_0$, since $\frac{\partial F}{\partial y}>0$, $(x,y)$ is strictly increasing as a function of $y$, therefore $! \exists y$ such that $F(x,y)=0$.
    This proves that $y=f(x)$ exists and is unique. Now,
    \begin{align*}
        F(x,f(x))&=0\\
        \frac{\partial F}{\partial x} \cdot 1+ \frac{\partial F}{\partial y} \frac{d f}{d x} &= 0
        \implies \frac{d f}{d x} = -\frac{\frac{\partial F}{\partial x}}{\frac{\partial F}{\partial y}}.
    \end{align*}
\end{proof}
\begin{theorem}
    $$
    \begin{cases}
        F_1(x_1,\cdots,x_n,z_1,\cdots,z_m)&=0\\
        \vdots\\
        F_m(x_1,\cdots,x_n,z_1,\cdots,z_m)&=0
    \end{cases}
    $$
    where $F_i\in C^1$ for $i=1,\cdots,m$. Let 
\[\triangle=\det\begin{pmatrix}
\frac{\partial F_1}{\partial z_1} & \cdots & \frac{\partial F_1}{\partial z_m}\\
\vdots & \ddots & \vdots\\
\frac{\partial F_m}{\partial z_1} & \cdots & \frac{\partial F_m}{\partial z_m}
\end{pmatrix}\]
If $\triangle\neq 0$ at a point $(\vec x^0,\vec z^0)$, then near that point, there are unique functions $z_i = f_i(x_1,\cdots,x_n)$ for $i=1,\cdots,m$ such that the system of equations is satisfied.
and the partial derivatives of $f_i$'s can be computed using implicit differentiation.
\end{theorem}
\begin{example}
Consider the system of equations:
\[
\begin{cases}
    xu+yvu^2&=2\\
    xu^3+y^2v^4&=2
\end{cases}
\]
Let's show that near $(x,y,u,v)=(1,1,1,1)$, we can express $u$ and $v$ as functions of $x$ and $y$.
\[
\begin{cases}
F_1(x,y,u,v) &= xu+yvu^2-2\\
F_2(x,y,u,v) &= xu^3+y^2v^4-2
\end{cases}
\]
Here, \[
\triangle = \det\begin{pmatrix}
\frac{\partial F_1}{\partial u} & \frac{\partial F_1}{\partial v}\\
\frac{\partial F_2}{\partial u} & \frac{\partial F_2}{\partial v}
\end{pmatrix} = \det\begin{pmatrix}x+2yvu & yu^2\\
3xu^2 & 4y^2v^3 \end{pmatrix}.
\]
Evaluating at the point $(1,1,1,1)$ gives:
\[\triangle(1,1,1,1) = \det\begin{pmatrix}1+2(1)(1) & 1(1)^2\\
3(1)(1)^2 & 4(1)^2(1)^3 \end{pmatrix} = \det\begin{pmatrix}3 & 1\\
3 & 4 \end{pmatrix} = 3\cdot4 - 1\cdot3 = 12 - 3 = 9 \neq 0.\]
By the implicit function theorem, there exist unique functions \(u=f_1(x,y)\) and \(v=f_2(x,y)\) near the point $(1,1,1,1)$.
Let's find $\frac{\partial u}{\partial x}(1,1)$:
\[
\begin{cases}
    u+x\frac{\partial u}{\partial x}+y\frac{\partial v}{\partial x}u^2+2yvu\frac{\partial u}{\partial x}&=0\\
    u^3+3xu^2\frac{\partial u}{\partial x}+y^2\cdot 4v^3\frac{\partial v}{\partial x}&=0
\end{cases} \implies 
\begin{cases}
    3\frac{\partial u}{\partial x}+\frac{\partial v}{\partial x}&=-1\\
    3\frac{\partial u}{\partial x}+4\frac{\partial v}{\partial x}&=-1
\end{cases}
\]  
Now, solve the system of equations to find $\frac{\partial u}{\partial x}(1,1)$.
\end{example}